\subsection{The first torsion level: a sharp tame ramification range}

\begin{theorem}[B-RD7: the first-level tame range]\label{thm:B-RD7}
\status{Proved; independently audited by the parent agent.}
Let $p\geq5$, and permit arbitrary perfect residue fields of characteristic
$p$. There exists a complete discretely valued $K$ of characteristic zero
and absolute ramification index $e$, and an elliptic curve $E/K$ with
$E[p]\subset E(K)$, additive potentially good reduction, and semistability
defect $m$, if and only if
\[
 m\in\{2,3,4,6\},\qquad m\mid(p+1),\qquad
 (p-1)\mid e,\qquad e\geq\frac{p^2-1}{m}.
\]
In the existence direction one may take the residue field to be
$\overline{\mathbf F}_p$. In particular the set of possible $e$ is
\[
 \left\{e\in(p-1)\mathbf Z_{>0}:
 e\geq\frac{p^2-1}{m_{\max}(p)}\right\},
\]
where $m_{\max}(p)=2,6,4,6$ for $p\equiv1,5,7,11\pmod {12}$,
respectively. For $p=3$, the same assertion with a specified tame defect
$m=2$ or $4$ holds: defect $2$ occurs exactly at even $e\geq4$, and
defect $4$ occurs at every positive even $e$. No assertion about wild
3-primary inertia is included in the latter statement.
\end{theorem}

\begin{proof}
We first prove necessity. Completion of the maximal unramified extension
does not change $e$ or the inertia image. The ordinary case is excluded
by Proposition~\ref{prop:B-RD3}. For $p\geq5$, potential-good inertia is
tame by \cite[Corollary~2 of Theorem~2]{SerreTate1968Good}. Thus the
minimal good-reduction extension $L/K$ is cyclic of order $m$, with
supersingular reduction $C$. Over an algebraically closed field of
characteristic greater than three, writing $C$ in short Weierstrass form
shows that its origin-preserving automorphisms have the form
$(x,y)\mapsto(u^2x,u^3y)$, with $u^4a=a$ and $u^6b=b$.
The possibilities are $m=2,3,4,6$. The curves with $j=0$ are supersingular
exactly when $p\equiv-1\pmod3$, and those with $j=1728$ exactly when
$p\equiv-1\pmod4$. Both assertions follow directly by extracting the
coefficient of $x^{p-1}$ from $(x^3+b)^{(p-1)/2}$ and
$(x^3+ax)^{(p-1)/2}$, respectively. Consequently $m\mid p+1$.
The same discussion for tame inertia in characteristic three gives
$m=2$ or $4$.

Normalize the valuation $v_L$ by $v_L(L^\times)=\mathbf Z$.
An average over the tame cyclic action provides a formal parameter $t$
at the identity of the good model which is an exact eigenvector.
If its eigenvalue is a primitive $m$th root $\zeta^q$, then every
$K$-rational nonzero $p$-torsion point $P$ satisfies
\[
 v_L(t(P))\equiv q\pmod m,\qquad \gcd(q,m)=1.
\]
All these valuations are positive, since $C$ is supersingular. Choose a
point $P$ whose parameter has maximum valuation $A$, and a point $Q$
whose parameter has minimum valuation $B$. If $A>B$ they are independent;
if $A=B$, choose any basis. The $p-1$ nonzero multiples of $P$ have
valuation $A$. The other $p(p-1)$ points have valuation $B$: if $A>B$,
this follows from the linear term of the formal group law; if $A=B$,
it follows from the assumed common maximum and minimum.
The constant term of $[p](t)/t$ is $p$, and its distinguished degree is
$p^2-1$. Weierstrass preparation therefore gives
\[
 m e=v_L(p)=(p-1)(A+pB).                         \tag{B.7.1}
\]
Here $A\geq B\geq1$ and $A\equiv B\equiv q\pmod m$.
Since $p\equiv-1\pmod m$, the displayed identity implies both
$(p-1)\mid e$ and $me\geq p^2-1$.

We give a deformation-theoretic converse with the equivariance made
explicit. Put $k=\overline{\mathbf F}_p$. Choose a supersingular elliptic
curve $C/k$ with an automorphism $\alpha$ of order $m$. For $m=3,6$
choose $j=0$, for $m=4$ choose $j=1728$ (the unique supersingular
$j$ in characteristic three is also zero), and for $m=2$ choose any
supersingular curve. The derivative of $\alpha$ is a primitive $m$th
root; a finite automorphism of order prime to $p$ with derivative one
is the identity, by averaging a local coordinate.

The framed deformation functor of $C$, with full Drinfeld $p$-level
structure, is represented by a complete regular local ring $R$ of
dimension two. If $P,Q$ are the universal labelled basis, their formal
coordinates form a regular system of parameters. These are the
precise statements of \cite[Theorem~2.1.2 and
Proposition~2.2.3(i)]{CarterStrauch2021LubinTate}; Serre--Tate identifies
the height-two formal-group deformation with the elliptic deformation
\cite[Theorem~1.2.1]{Katz1981SerreTate}. The universal formal elliptic
curve algebraizes: the identity gives the compatible ample line
bundles $\mathcal O(3[0])$, so the projective formal existence theorem
applies over the complete noetherian ring $R$.

Changing the framing by $\alpha$ gives an action of $C_m$ on this
functor, hence a semilinear action on the universal elliptic curve.
It preserves the labels $P,Q$. For any initial formal coordinate $t_0$,
replace it by
\[
 t=\frac1m\sum_{i=0}^{m-1}\zeta^{-i}\alpha^{i*}t_0,
\]
where $\zeta$ is the derivative eigenvalue. The linear coefficient is
a unit, so $t$ remains a coordinate and $\alpha^*t=\zeta t$.
Thus $x=t(P)$ and $y=t(Q)$ are regular parameters satisfying exactly
$\alpha(x)=\zeta x$, $\alpha(y)=\zeta y$. The action fixes $W(k)$.
Cohen's presentation theorem gives an equivariant surjection
\[
 W(k)[[X,Y]]\longrightarrow R,\qquad X\mapsto x,\quad Y\mapsto y,
\]
where $\alpha$ multiplies $X,Y$ by $\zeta$.

We detail its single relation, to avoid assuming that an arbitrary
monomial specialization has the intended ramification. The full-level
factorization gives, in $R$,
\[
 p=u\prod_{(a,b)\in\mathbf F_p^2\setminus\{0\}}
          t([a]P+[b]Q),\qquad u\in R^\times.
\]
Lift $u$ to a unit $U$ and each factor to a power series $T_{a,b}$
in $W(k)[[X,Y]]$ by using the multiplication and addition power
series of the formal group, lifting their coefficients. These lifts
have zero constant term and linear term $aX+bY$ (using integral
representatives for $a,b$); crucially, $T_{a,0}$ is divisible by $X$
with unit quotient when $a\ne0$. Then
\[
 F_0=p-U\prod_{(a,b)\ne0}T_{a,b}
\]
lies in the kernel and generates it: it has coefficient one on $p$
modulo the square of the maximal ideal of $W(k)[[X,Y]]$.
Replace $F_0$ by its $C_m$-average $F$. This still generates the kernel,
has the same linear term $p$, and is invariant.

Write $t=e/(p-1)$ and choose
\[
 B=1,\qquad A=1+m\left(t-\frac{p+1}{m}\right).
\]
These are positive integers, $A\geq B$, and $A\equiv B\equiv1\pmod m$.
Reduce the relation modulo $p$, and give $X,Y$ weights $A,B$. If $A>B$, the weighted initial term of the
product in $F_0$ is a nonzero scalar times
$X^{p-1}Y^{p(p-1)}$, of weight
$d=(p-1)(A+pB)$. If $A=B$, the weighted initial term is the nonzero
homogeneous polynomial $\prod_{(a,b)\ne0}(aX+bY)$ of degree $p^2-1$.
In either case that initial form is invariant, since its total degree
is $p^2-1$, divisible by $m$. Averaging therefore preserves it.
Choose a Teichmuller unit $c$ arbitrarily when $A>B$, and with
$\bar c\notin\mathbf F_p$ when $A=B$. Then
\[
 R_L=W(k)[[\pi]]/(F(\pi^A,c\pi^B))
\]
is a complete discrete valuation ring with uniformizer $\pi$ and
$v_L(p)=d$: modulo $p$ its defining equation has first nonzero term
of degree $d$, while its coefficient on $p$ is a unit. This also
proves directly that the specified specialization exists.

The action $\pi\mapsto\zeta\pi$ preserves this equation, by invariance
of $F$ and $A\equiv B\equiv1\pmod m$. It has order $m$, acts trivially
on the residue field, and makes the map $R\to R_L$ equivariant.
Let $L=\operatorname{Frac}R_L$ and $K=L^{C_m}$. Then $L/K$ is tame
of ramification degree $m$, and
\[
 e(K)=d/m=e.
\]
The generic elliptic curve descends to $K$. Since the labels of the
level structure were preserved, its full $p$-torsion is rational over
$K$. The inertia action on the good special fibre is the faithful
order-$m$ action $\alpha$. The Neron--Ogg--Shafarevich criterion gives
exactly semistability defect $m$, so reduction is additive. Finally the
Weil pairing implies $\mu_p\subset K$. This proves every sufficiency
assertion made in the theorem, including $p=3$ and $m=2,4$.
\end{proof}
